% !TEX root = ../../main.tex
% 来源：《代数学引论》第三版，第四章，第137--157页。
\setchapterimage{}
\setchapterintro{整环中的整除、主理想、唯一因子分解，以及多项式扩张与诺特性。}
\chapter{整环的整除性}\label{p2:ch4}

\section{主理想整环}

本章的整环均为交换幺环，且$1\ne0$。记$R$的单位群为$U$。

\begin{definition}
对$a,b\in R$，若存在$c\in R$使$a=bc$，称$b$整除$a$，记为$b\mid a$；称$b$为$a$的因子，$a$为$b$的倍数。

若$a\mid b$且$b\mid a$，称$a,b$相伴，记为$a\sim b$。若$b\mid a$而$a\nmid b$，称$b$为$a$的真因子。

非零非单位元素$a$若由$a=bc$必有$b\sim1$或$c\sim1$，称$a$为不可约元。非零非单位元素$a$若由$a\mid bc$必有$a\mid b$或$a\mid c$，称$a$为素元。
\end{definition}

整除具有自反性、传递性，且$c\mid a,c\mid b$蕴含$c\mid(au+bv)$，$u,v\in R$。$1$的全部因子恰为单位。相伴是等价关系，而且
\[
a\sim b\quad\Longleftrightarrow\quad b=ua\text{ 对某个 }u\in U\text{ 成立}.
\]
事实上，若$a\ne0$、$b=va$、$a=wb$，消去$a$得$wv=1$；$a=0$时，相伴迫使$b=0$。因此$a$的相伴类为$Ua$，非零元的因子按相伴分类，$U$和$Ua$中的因子为平凡因子。

\begin{definition}
若$d\mid a,d\mid b$，且任一公因子$c$均满足$c\mid d$，则称$d$为$a,b$的最大公因子。若$a\mid m,b\mid m$，且对任一公倍数$c$都有$m\mid c$，则称$m$为最小公倍数。它们若存在，都只在相伴意义下唯一。
\end{definition}

\begin{lemma}\label{p2:divisibility-ideals}
设$R$为整环，则：
\begin{enumerate}
\item $a\mid b$当且仅当$(b)\subseteq(a)$；$a\sim b$当且仅当$(a)=(b)$；真因子对应严格包含。
\item $a$为素元当且仅当$(a)$为非零素理想；若$(a)$为非零极大理想，则$a$为素元。
\item 素元一定是不可约元。
\end{enumerate}
\end{lemma}
\begin{proof}
\begin{enumerate}
\item 若$b=ac$，则$Rb\subseteq Ra$。反过来，$(b)\subseteq(a)$给出$b\in(a)$，所以$b=ac$。双向应用此结论得到相伴与严格包含的刻画。
\item $bc\in(a)$等价于$a\mid bc$，而$b\in(a)$等价于$a\mid b$，故两种素性的定义等价。非零素理想的生成元既非零又非单位。极大理想为素理想，所以最后的结论也成立。
\item 若素元$a=bc$，则$a\mid b$或$a\mid c$。若$b=ad$，代入$a=adc$并消去$a\ne0$，得$dc=1$，故$c$为单位；另一情形同理。因此$a$不可约。
\end{enumerate}
\end{proof}

\begin{theorem}\label{p2:pid-primes}
在主理想整环$R$中，不可约元都是素元；每个非零素理想都是极大理想。
\end{theorem}
\begin{proof}
先设$a$不可约。若$(a)\subseteq N\subseteq R$，写$N=(b)$，则$b\mid a$。由不可约性，$b\sim a$或$b\sim1$，故$N=(a)$或$N=R$。因此$(a)$为极大理想，进而$a$为素元。

再设$N$为非零素理想。写$N=(a)$，由引理\ref{p2:divisibility-ideals}，$a$为素元，故不可约。由刚才的证明，$N=(a)$极大。
\end{proof}

\begin{theorem}\label{p2:pid-bezout}
若$R$为主理想整环，且$(a)+(b)=(d)$，则$d$为$a,b$的最大公因子，而且存在$u,v\in R$使
\[
d=ua+vb.
\]
进一步，若$(a_1)+\cdots+(a_r)=(d)$，则$d$为这些元素的最大公因子，并可写成$\sum_{i=1}^r u_i a_i$。
\end{theorem}
\begin{proof}
由$(a),(b)\subseteq(d)$可知$d\mid a,d\mid b$。若$c\mid a,c\mid b$，则$(a),(b)\subseteq(c)$，故$(d)\subseteq(c)$，即$c\mid d$。而$d\in(a)+(b)$恰给出所列线性组合。有限多个元素的情形用同样的包含关系即可证明。
\end{proof}

\begin{definition}
理想列$N_1\subseteq N_2\subseteq\cdots$称为理想升链。若元素列满足$a_{i+1}\mid a_i$，称为因子降链；若从某项起$a_m\sim a_{m+1}\sim\cdots$，称此降链终止。理想链从某项起相等，也称终止。
\end{definition}

\begin{lemma}\label{p2:pid-acc}
设$R$为主理想整环，则：
\begin{enumerate}
\item $R$的任一理想升链$\{(a_i)\}$恒有一个正整数$m$，使得
\[
(a_m)=(a_{m+1})=(a_{m+2})=\cdots.
\]
\item $R$的任一因子降链$\{a_i\}$恒有一个正整数$m$，使得
\[
a_m\sim a_{m+1}\sim a_{m+2}\sim\cdots.
\]
\end{enumerate}
\end{lemma}
\begin{proof}
根据引理\ref{p2:divisibility-ideals}，
\[
a_{i+1}\mid a_i
\quad\Longleftrightarrow\quad
(a_i)\subseteq(a_{i+1}),
\qquad
 a_i\sim a_{i+1}
\quad\Longleftrightarrow\quad
(a_i)=(a_{i+1}).
\]
又因为$R$的每个理想都是主理想，所以两条结论等价。现证明第一条。

设$(a_1)\subseteq(a_2)\subseteq\cdots$为$R$的一条理想升链，令
\[
N=\bigcup\limits_{i\ge1}(a_i).
\]
首先证明$N$是一个理想。每个$(a_i)$都含有$0$，所以$N$非空。对$a,b\in N$，存在正整数$r,s$使
\[
a\in(a_r),\qquad b\in(a_s).
\]
不妨设$r\le s$，于是$(a_r)\subseteq(a_s)$，从而$a,b\in(a_s)$。因为$(a_s)$是理想，所以
\[
a-b\in(a_s)\subseteq N.
\]
对$c\in R$，由$a\in(a_r)$得
\[
ac\in(a_r)\subseteq N.
\]
由于$R$交换，也有$ca=ac\in N$。因此$N$是$R$的一个理想。

因为$R$是主理想整环，所以存在$d\in R$使得$N=(d)$。根据$N$的定义，$d$属于某个$(a_m)$。由于$(a_m)$为包含$d$的理想，便有
\[
N=(d)\subseteq(a_m).
\]
反之，由并集的定义有$(a_m)\subseteq N$，所以
\[
N=(a_m).
\]
对任意大于$m$的整数$n$，由理想升链的包含关系及$N$的定义，
\[
(a_m)\subseteq(a_n)\subseteq N.
\]
结合$(a_m)=N$，得到$(a_n)=N$。最后得到
\[
N=(a_m)=(a_{m+1})=(a_{m+2})=\cdots.
\]
第一条得证，第二条也随之成立。
\end{proof}

\begin{example}
$\Z[x]$不是主理想整环。考虑$I=(2,x^2+1)$，若$I=(g)$，由$2\in(g)$及次数公式，$g$为整常数，且$g\mid2$，所以忽略符号后只能$g=1$或$2$。

若$g=1$，则$1=2u(x)+(x^2+1)v(x)$。代入$x=1$得$1=2u(1)+2v(1)$，矛盾。若$g=2$，则$x^2+1$的全部系数均为偶数，也矛盾。所以$I$不是主理想。
\end{example}

\section{欧几里得整环}

\begin{definition}
设$R$为整环。若存在函数$d:R\setminus\{0\}\to\N$，使对任意$a,b\in R$、$b\ne0$，存在$q,r\in R$满足
\[
a=qb+r,\qquad r=0\text{ 或 }d(r)<d(b),
\]
则称$R$为欧几里得整环。
\end{definition}

例如$\Z$可取$d(a)=|a|$，域$F$上的$F[x]$可取$d(f)=1+\deg f$。定义不要求商和余数唯一，也不要求$d$是唯一选定的函数。

\begin{theorem}\label{p2:euclidean-pid}
欧几里得整环是主理想整环。
\end{theorem}
\begin{proof}
设$N$为理想，$N=(0)$时已为主理想。若$N\ne(0)$，从其非零元素中选$b$，使$d(b)$最小。对任意$a\in N$作除法，得$a=qb+r$，其中$r=0$或$d(r)<d(b)$。

因$r=a-qb\in N$，若$r\ne0$就违反$b$的选择，故$r=0$，即$a\in(b)$。于是$N\subseteq(b)$；由$b\in N$又有$(b)\subseteq N$，所以$N=(b)$。
\end{proof}

在欧几里得整环中可逐次作带余除法。非零余数的$d$值严格下降，因此算法有限步终止；利用$(a,b)=(b,r)$所对应的公因子相同关系，最后一个非零余数为最大公因子，回代可得到它的线性组合表示。

\begin{definition}
设$m\ne0,1$为无平方因子的整数，即不存在素数$p$使$p^2\mid m$。称
\[
\Q(\sqrt m)=\{s+t\sqrt m:s,t\in\Q\}
\]
为一个\textbf{二次域}。由于$\sqrt m\notin\Q$，每个元素的表示$s+t\sqrt m$唯一。

对$\alpha=s+t\sqrt m$，称$\overline\alpha=s-t\sqrt m$为$\alpha$的\textbf{共轭}，称
\[
N(\alpha)=\alpha\overline\alpha=s^2-mt^2
\]
为$\alpha$的\textbf{范数}。若$\alpha$在$\Q$上的首一极小多项式的系数均为整数，则称$\alpha$为\textbf{代数整数}。
\end{definition}

\begin{proposition}
$\Q(\sqrt m)$的代数整数全体为
\[
R_m=\begin{cases}
\Z[\sqrt m],&m\equiv2,3\pmod4,\\
\Z[(1+\sqrt m)/2],&m\equiv1\pmod4.
\end{cases}
\]
它是$\Q(\sqrt m)$的子环，称为该二次域的\textbf{代数整数环}。
\end{proposition}
\begin{proof}
设$\alpha=s+t\sqrt m$，其中$s,t\in\Q$。

先考虑$t=0$的情形。此时$\alpha=s\in\Q$，其在$\Q$上的首一极小多项式为$x-s$。它的系数均为整数当且仅当$s\in\Z$，所以有理数中的代数整数恰为整数。

以下设$t\ne0$。由$\sqrt m\notin\Q$可知$\alpha\notin\Q$，所以$\alpha$不可能满足有理系数的一次多项式。另一方面，
\begin{align*}
(x-\alpha)(x-\overline\alpha)
&=x^2-(\alpha+\overline\alpha)x+\alpha\overline\alpha\\
&=x^2-2sx+s^2-mt^2
\end{align*}
是以$\alpha$为根的首一有理系数二次多项式。因此它就是$\alpha$在$\Q$上的首一极小多项式。于是
\[
\alpha\text{为代数整数}
\quad\Longleftrightarrow\quad
2s\in\Z\text{ 且 }s^2-mt^2\in\Z.
\]

设这两个条件成立，记$a=2s\in\Z$、$k=s^2-mt^2\in\Z$。将$t$写成最简分数$t=u/v$，其中$u,v\in\Z$、$v>0$且$\gcd(u,v)=1$。由$4mt^2=a^2-4k$得
\[
4mu^2=v^2(a^2-4k).
\]
所以$v^2\mid4mu^2$。由于$\gcd(u^2,v^2)=1$，可以消去与$v^2$互素的因子$u^2$，得到
\[
v^2\mid4m.
\]
若奇素数$p$整除$v$，则$p^2\mid4m$，从而$p^2\mid m$，与$m$无平方因子矛盾。因此$v$只能是$2$的幂。若$4\mid v$，则$16\mid v^2$，进而$16\mid4m$，即$4\mid m$，仍与$m$无平方因子矛盾。故$v=1$或$v=2$，于是
\[
b=2t\in\Z,\qquad
\alpha=\frac{a+b\sqrt m}{2}.
\]
此时$s^2-mt^2=(a^2-mb^2)/4$，所以必须有
\[
a^2-mb^2\equiv0\pmod4.
\]
反过来，若$a,b\in\Z$满足这一同余式，并令$s=a/2$、$t=b/2\ne0$，则$2s=a\in\Z$且$s^2-mt^2\in\Z$，故相应的$\alpha$为代数整数。

$t=0$时，上述表示中的$b=0$，同余条件为$a^2\equiv0\pmod4$，等价于$a$为偶数，恰好给出$\alpha=a/2\in\Z$。因此，在所有情形下均有
\[
R_m=\left\{\frac{a+b\sqrt m}{2}:a,b\in\Z,\ a^2-mb^2\equiv0\pmod4\right\}.
\]

因为$m$无平方因子，所以$4\nmid m$，只需讨论$m$模$4$余$1,2,3$的情形。偶数的平方模$4$为$0$，奇数的平方模$4$为$1$。
\begin{enumerate}
\item 当$m\equiv2,3\pmod4$时，若$b$为奇数，则$b^2\equiv1\pmod4$，于是
\[
a^2\equiv mb^2\equiv m\pmod4,
\]
即$a^2$模$4$余$2$或$3$，不可能。因此$b$为偶数，从而$mb^2\equiv0\pmod4$，同余条件又迫使$a$为偶数。

反过来，若$a,b$均为偶数，则$a^2-mb^2$必被$4$整除。所以此时
\[
R_m=\left\{\frac{a+b\sqrt m}{2}:a,b\in2\Z\right\}
=\{c+d\sqrt m:c,d\in\Z\}.
\]
由于$(\sqrt m)^2=m\in\Z$，$\sqrt m$的任意整数系数多项式都可以化为$c+d\sqrt m$，而这种形式的元素也都属于$\Z[\sqrt m]$。故$R_m=\Z[\sqrt m]$。

\item 当$m\equiv1\pmod4$时，同余条件等价于
\[
a^2\equiv b^2\pmod4.
\]
由整数平方模$4$的取值可知，该条件等价于$a,b$同为偶数或同为奇数。

令$\omega=(1+\sqrt m)/2$。当$a,b$同奇偶时，$(a-b)/2\in\Z$，且
\[
\frac{a+b\sqrt m}{2}=\frac{a-b}{2}+b\omega.
\]
反过来，对任意$c,d\in\Z$，有
\[
c+d\omega=\frac{(2c+d)+d\sqrt m}{2},
\]
其中$2c+d$与$d$同奇偶，所以$c+d\omega$为代数整数。因此$R_m=\{c+d\omega:c,d\in\Z\}$。

又因为$h=(m-1)/4\in\Z$，并且
\[
\omega^2=\frac{1+2\sqrt m+m}{4}=\omega+h,
\]
利用这个等式，$\omega$的任意整数系数多项式都可化为$c+d\omega$。因此$R_m=\Z[\omega]$。
\end{enumerate}

最后直接验证上述集合为子环。$m\equiv2,3\pmod4$时，对$c,d,e,f\in\Z$，
\begin{align*}
(c+d\sqrt m)-(e+f\sqrt m)
&=(c-e)+(d-f)\sqrt m,\\
(c+d\sqrt m)(e+f\sqrt m)
&=(ce+mdf)+(cf+de)\sqrt m.
\end{align*}
两式中的系数均为整数，因此$\Z[\sqrt m]$对减法与乘法封闭，并且含有$0$和$1$，故为子环。

$m\equiv1\pmod4$时，仍记$h=(m-1)/4\in\Z$。由$\omega^2=\omega+h$，
\begin{align*}
(c+d\omega)-(e+f\omega)
&=(c-e)+(d-f)\omega,\\
(c+d\omega)(e+f\omega)
&=ce+(cf+de)\omega+df\omega^2\\
&=(ce+hdf)+(cf+de+df)\omega.
\end{align*}
两式仍为$\omega$的整数线性组合，并且$\Z[\omega]$含有$0$和$1$，故也为子环。
\end{proof}

\begin{example}
高斯整数环$R_{-1}=\Z[\sqrt{-1}]$是欧几里得整环，取$d(\alpha)=N(\alpha)$。对$\beta\ne0$，写$\alpha/\beta=t+s\sqrt{-1}$，选整数$u,v$使$|t-u|,|s-v|\le1/2$，令
\[
q=u+v\sqrt{-1},\qquad r=\alpha-q\beta.
\]
则$r\in R_{-1}$，且范数的乘法性给出
\[
N(r)=N(\beta)\bigl((t-u)^2+(s-v)^2\bigr)
\le\tfrac12N(\beta)<N(\beta).
\]
所以满足欧几里得除法条件。
\end{example}

\begin{example}
$R_{-3}$也是欧几里得整环，仍取范数。写$\alpha/\beta=t+s\sqrt{-3}$。先选整数$v$使$|2s-v|\le1/2$，再选与$v$同奇偶的整数$u$使$|2t-u|\le1$。后一种整数相邻间距为$2$，所以这样的$u$存在。令
\[
q=\tfrac12(u+v\sqrt{-3}),\qquad r=\alpha-q\beta.
\]
因为$u,v$同奇偶，$q\in R_{-3}$，从而$r\in R_{-3}$。而
\[
N(r)=N(\beta)\left[\left(\frac{2t-u}{2}\right)^2
+3\left(\frac{2s-v}{2}\right)^2\right]
\le\frac7{16}N(\beta)<N(\beta).
\]
故除法条件成立。
\end{example}

% 原文另引用二次域分类结果，不在本章证明其分类。
对于虚二次域，代数整数环为主理想整环的$m$恰为
\[
-1,-2,-3,-7,-11,-19,-43,-67,-163.
\]
前五个可用范数给出欧几里得除法，后四个虽为主理想整环，却不是欧几里得整环。对于实二次域，以$|N(\alpha)|$为欧几里得函数的$m$恰为
\[
2,3,5,6,7,11,13,17,19,21,29,33,37,41,57,73.
\]


\section{唯一因子分解整环}

\begin{definition}
若整环$R$满足下列条件，称它为唯一因子分解整环，也称高斯整环：
\begin{enumerate}
\item 每个非零非单位元可写成有限个不可约元之积；
\item 若$a=p_1\cdots p_r=q_1\cdots q_s$为两个这样的分解，则$r=s$，而且适当调整次序后$q_i\sim p_i$。
\end{enumerate}
若每条因子降链均终止，称$R$满足因子降链条件。
\end{definition}

\begin{lemma}\label{p2:acc-factorization}
若整环满足因子降链条件，则每个非零非单位元都可分解为有限个不可约元之积。
\end{lemma}
\begin{proof}
反设不能分解的非零非单位元组成非空集合$S$。对$a\in S$，$a$不是不可约元，所以$a=bc$，其中$b,c$均非单位。若$b,c$都可分解，则$a$也可分解，矛盾。因此至少一个真因子仍在$S$中。

从$a_0\in S$出发，逐次选取$a_{n+1}\in S$为$a_n$的真因子，便得
\[
a_{n+1}\mid a_n,\qquad a_n\nmid a_{n+1}.
\]
这是一条不终止的因子降链，与假设矛盾。
\end{proof}

\begin{theorem}\label{p2:ufd-criterion}
若整环$R$满足因子降链条件，且每个不可约元均为素元，则$R$为唯一因子分解整环。特别地，主理想整环是唯一因子分解整环。
\end{theorem}
\begin{proof}
存在性由前一引理得到。对唯一性，设$p_1\cdots p_r=q_1\cdots q_s$。因$p_1$为素元，它整除某个$q_j$；重新排列使$j=1$。$q_1$不可约，而$p_1$为非单位，故$q_1=u p_1$，$u$为单位。消去$p_1$，得到
\[
p_2\cdots p_r=u q_2\cdots q_s.
\]
重复此过程。若一边先无因子，另一边仍有不可约因子，则一个单位等于非单位元素的乘积，不可能。因此$r=s$，且对应因子相伴。

主理想整环满足降链条件，且不可约元为素元，所以也满足结论。
\end{proof}

\begin{lemma}\label{p2:gcd-scaling}
若$R$中每对元素均有最大公因子，则对$c\ne0$，$c(a,b)$是$ca,cb$的最大公因子。
\end{lemma}
\begin{proof}
设$d=(a,b)$、$e=(ca,cb)$。$cd$为$ca,cb$的公因子，故$e=cdu$。另一方面，$c$为公因子，故可写$e=ce_1$。由$e\mid ca,cb$，消去$c$得$e_1\mid a,b$，所以$e_1\mid d$。又$e_1=du$，故$d=duv$对某个$v$成立。

若$d\ne0$，消去$d$得$uv=1$，所以$e\sim cd$，证明结论。$d=0$时$a=b=0$，结论直接成立。
\end{proof}

\begin{lemma}
若整环中每对元素均有最大公因子，则不可约元都是素元。
\end{lemma}
\begin{proof}
设$p$不可约，且$p\mid ab$。令$d=(p,a)$，则$d\sim p$或$d\sim1$。前者给出$p\mid a$。后者把$d$取为$1$，由引理\ref{p2:gcd-scaling}，$b$为$pb,ab$的最大公因子。由于$p$同时整除$pb$和$ab$，得到$p\mid b$。故$p$为素元。
\end{proof}

\begin{corollary}
满足因子降链条件、且每对元素都有最大公因子的整环，是唯一因子分解整环。
\end{corollary}

\begin{proposition}
在唯一因子分解整环中，从不可约元的每个相伴类选一个代表$p\in S$，则每个非零元素都有唯一的标准分解
\[
a=u\prod_{p\in S}p^{r_p},\qquad u\in U,quad r_p\ge0,
\]
其中只有有限个$r_p$非零。若$b=v\prod p^{s_p}$，则：
\begin{enumerate}
\item $a\mid b$当且仅当$r_p\le s_p$对每个$p$成立；$a\sim b$当且仅当全部指数相同；
\item $a$的因子按相伴分类共有$\prod_p(r_p+1)$类，且因子降链条件成立；
\item $\prod_p p^{\min(r_p,s_p)}$为最大公因子，$\prod_p p^{\max(r_p,s_p)}$为最小公倍数。
\end{enumerate}
\end{proposition}
\begin{proof}
\begin{enumerate}
\item 若$b=ac$，将$c$也分解，唯一性说明$b$中各指数为$a,c$中相应指数之和。反之，指数差均非负时直接作乘积即可构造$c$。
\item 因子中$p$的指数可独立选为$0,\ldots,r_p$，得到计数。对以$a$开始的严格因子降链，每下降一步指数总和至少减$1$，故只能有限步。
\item 同时整除$a,b$等价于每个指数不超过两者的最小值；同时为它们的倍数等价于每个指数至少为两者的最大值，因而得到两个结论。
\end{enumerate}
\end{proof}

\begin{theorem}\label{p2:ufd-pid-characterization}
设$R$为唯一因子分解整环，则下列条件均等价于$R$为主理想整环：
\begin{enumerate}
\item 任意$a,b$的一个最大公因子可写为$ua+vb$；
\item 每个不可约元生成的理想都是极大理想。
\end{enumerate}
\end{theorem}
\begin{proof}
\begin{enumerate}
\item 必要性由定理\ref{p2:pid-bezout}得到。证充分性。设$N\ne(0)$为理想，取$0\ne a_1\in N$。若$N\ne(a_1)$，取$b_1\in N\setminus(a_1)$。由假设，$(a_1)+(b_1)=(a_2)$，其中$a_2$为最大公因子，故
\[
(a_1)\subsetneq(a_2)\subseteq N.
\]
若$N\ne(a_2)$，重复此构造。严格包含对应非相伴真因子，因子降链条件保证过程终止，最终$N=(a_n)$。零理想也为主理想。
\item 必要性已证。反过来，不相伴的不可约元$a,b$生成不同的极大理想，故$(a)+(b)=R$。对任意非零$a,b$，写$a=da_1,b=db_1$，其中$d=(a,b)$。$a_1,b_1$无共同不可约因子，因此任一素因子对生成的理想都互素。反复使用第三章中“分别与一个理想互素的两个理想的积仍与它互素”，得到$(a_1)+(b_1)=R$。于是
\[
(a)+(b)=(d)\bigl((a_1)+(b_1)\bigr)=(d).
\]
若某个$a_1,b_1$为单位，此等式同样成立。零元情形直接处理。因而最大公因子为线性组合，再由第一条，$R$为主理想整环。
\end{enumerate}
\end{proof}

\begin{example}
$R=\Z[\sqrt{-5}]$不是唯一因子分解整环。范数$N(a+b\sqrt{-5})=a^2+5b^2$具有乘法性，故单位仅为$\pm1$。

考虑
\[
6=2\cdot3=(1+\sqrt{-5})(1-\sqrt{-5}).
\]
范数为$2$或$3$的元素不存在：若$b\ne0$则范数至少为$5$，若$b=0$则范数为平方。若$1+\sqrt{-5}$分解为两个非单位，其范数$6$必须分解为$2\cdot3$，不可能，故它不可约。共轭元素同理。$2$若可约，其范数$4$须分为$2\cdot2$；$3$若可约，其范数$9$须分为$3\cdot3$，也都不可能。

这四个不可约元两两不相伴：相伴只能相差符号，直接比较有理部和$\sqrt{-5}$的系数即可排除。因此两个分解不等价。
\end{example}

% 原文还说明理想分解如何替代元素分解，以下保留其具体构造。
令
\[
P_1=(2,1+\sqrt{-5}),\ P_2=(3,1+\sqrt{-5}),\quad
P_3=(2,1-\sqrt{-5}),\ P_4=(3,1-\sqrt{-5}).
\]
它们的商环分别为$\F_2,\F_3$，所以均为极大理想；由生成元的乘积可验证
\[
(2)=P_1P_3,\quad(3)=P_2P_4,\quad
(1+\sqrt{-5})=P_1P_2,\quad(1-\sqrt{-5})=P_3P_4.
\]
例如$P_1P_3=(4,2(1+\sqrt{-5}),2(1-\sqrt{-5}),6)$，全部生成元均属于$(2)$，且$6-4=2$在乘积理想中，故等于$(2)$。另外，$P_2P_4=(9,3(1+\sqrt{-5}),3(1-\sqrt{-5}),6)$，其中$9-6=3$；$P_1P_2$的生成元为$6,2(1+\sqrt{-5}),3(1+\sqrt{-5}),(1+\sqrt{-5})^2$，它们均为$1+\sqrt{-5}$的倍数，且后两个一次倍数相减得到$1+\sqrt{-5}$。$P_3P_4$对共轭元素作相同计算即可。因此元素的两种分解在理想层面给出同一个分解$(6)=P_1P_2P_3P_4$。

\section{高斯整环的多项式扩张}

\begin{definition}
设$R$为高斯整环。非零多项式$f=\sum_{i=0}^na_ix^i$的系数最大公因子称为$f$的容度，记为$c(f)$，只在相伴意义下确定。若$c(f)\sim1$，称$f$为本原多项式。
\end{definition}

\begin{lemma}\label{p2:content-decomposition}
每个非零$f\in R[x]$可写成$f=df_1$，其中$d\in R$，$f_1$为本原多项式，而且$d,f_1$均在相伴意义下唯一。
\end{lemma}
\begin{proof}
取$d=c(f)$，将每个系数写成$a_i=da_i'$，令$f_1=\sum a_i'x^i$。若$a_i'$有非单位公因子，则$da_i'$的最大公因子不可能仅与$d$相伴，所以$f_1$本原。

若另有$f=ef_2$，其中$f_2$本原，由最大公因子在数乘下的性质，$c(f)\sim d\sim e$。写$e=ud$，$u$为单位，消去$d$得到$f_1=uf_2$，故也相伴。
\end{proof}

\begin{lemma}[高斯引理]\label{p2:gauss-primitive}
两个本原多项式的积仍为本原多项式。
\end{lemma}
\begin{proof}
设$f=\sum a_ix^i$、$g=\sum b_jx^j$均本原，反设$h=fg$不本原。取不可约元$p$整除$c(h)$。因为本原性，存在最小的$r,s$使$p\nmid a_r$、$p\nmid b_s$。

在$h$的$x^{r+s}$系数中，$a_rb_s$之外的每项$a_ib_j$都被$p$整除：若$i<r$则$p\mid a_i$，若$i>r$则$j<s$，故$p\mid b_j$。而$p$为素元，$p\nmid a_rb_s$。因此$p$不能整除这一系数，与$p\mid c(h)$矛盾。所以$h$本原。
\end{proof}

\begin{lemma}\label{p2:primitive-fraction-field}
设$F$为$R$的商域，则：
\begin{enumerate}
\item $R[x]$的两个本原多项式在$R[x]$中相伴，当且仅当在$F[x]$中相伴；
\item $F[x]$的每个非零多项式可写成$(d/e)g$，$d,e\in R\setminus\{0\}$、$g\in R[x]$本原，且$g$在$R[x]$中相伴意义下唯一。
\end{enumerate}
\end{lemma}
\begin{proof}
\begin{enumerate}
\item 正向成立。反过来，若$f=(a/b)g$，则$bf=ag$。由本原性，等式两侧容度分别与$b,a$相伴，所以$a=ub$，$u\in U$，得到$f=ug$。
\item 把系数写为有限个分数，取非零公分母$e$，则$ef\in R[x]$。由引理\ref{p2:content-decomposition}，$ef=dg$，$g$本原，所以$f=(d/e)g$。若另有$f=(a/b)h$且$h$本原，则$g,h$在$F[x]$中相伴，由第一条，在$R[x]$中也相伴。
\end{enumerate}
\end{proof}

\begin{corollary}
正次数本原多项式$g\in R[x]$在$R[x]$中不可约，当且仅当它在$F[x]$中不可约。
\end{corollary}
\begin{proof}
若$g=f_1f_2$在$F[x]$中可约，两个因子次数均为正。写$f_i=\alpha_i g_i$，$g_i\in R[x]$本原，则$g=\alpha_1\alpha_2g_1g_2$。高斯引理说明$g_1g_2$本原，故由上一引理，$g$与$g_1g_2$在$R[x]$中相伴，得到正次数分解。

反过来，若本原$g$在$R[x]$中可约，则非常数因子的分解已经是$F[x]$中的分解；不可能有非单位常数因子，因为那将整除$g$的全部系数，违反本原性。所以两个因子次数都为正。
\end{proof}

\begin{theorem}\label{p2:polynomial-ufd}
高斯整环$R$上的一元多项式环$R[x]$仍为高斯整环。进一步，有限多个未定元的$R[x_1,\ldots,x_n]$也是高斯整环。
\end{theorem}
\begin{proof}
先证存在性。对非零非单位$f$，写$f=dg$，$g$本原。将$d$在$R$中分解为不可约元之积。对正次数的$g$按次数归纳：若它可约，则分解为两个次数更低的正次数多项式；由高斯引理及容度分解，可取这些因子本原。逐次分解得到
\[
f=p_1\cdots p_tq_1\cdots q_r,
\]
其中$p_i$为$R$中的不可约元，$q_j$为正次数本原不可约多项式。常数$p_i$在$R[x]$中仍不可约，因为常数的任何非零分解都只能有常数因子。

再证唯一性。若另有同类分解$f=p_1'\cdots p_{t'}'q_1'\cdots q_{r'}'$，由容度分解及高斯引理，
\[
p_1\cdots p_t\sim p_1'\cdots p_{t'}',\qquad
q_1\cdots q_r\sim q_1'\cdots q_{r'}'.
\]
第一式在$R$中给出$t=t'$及对应常数因子相伴。第二式放入$F[x]$，所有$q_j$在其中仍不可约；由域上多项式的唯一分解，$r=r'$，适当排列后$q_j,q_j'$在$F[x]$中相伴。它们本原，故在$R[x]$中也相伴，唯一性得证。

多元情形对变量个数归纳，逐次应用一元结论。
\end{proof}

\begin{theorem}[艾森斯坦判别法]\label{p2:eisenstein}
设$R$为高斯整环，$F$为其商域，$f=a_0+\cdots+a_nx^n$，$n>1$。若存在不可约元$p$满足
\[
p\mid a_i\ (0\le i<n),\qquad p^2\nmid a_0,\qquad p\nmid a_n,
\]
则$f$在$F[x]$中不可约，即在$R[x]$中不能分解为两个正次数多项式之积。
\end{theorem}
\begin{proof}
先反设$f=gh$为$R[x]$中的正次数分解，写$g=\sum_{i=0}^rb_ix^i$、$h=\sum_{j=0}^sc_jx^j$，$r,s>0$。由$a_0=b_0c_0$及$p\mid a_0$、$p^2\nmid a_0$，两个常数项中恰有一个被$p$整除，不妨$p\mid b_0,p\nmid c_0$。

由$p\nmid a_n=b_rc_s$可知$p\nmid b_r,c_s$。取最小$j\ge1$使$p\nmid b_j$，则$j\le r<n$，且
\[
a_j=b_0c_j+b_1c_{j-1}+\cdots+b_jc_0.
\]
前$j$项被$p$整除，末项不被$p$整除，故$p\nmid a_j$，矛盾。

若$f$不本原，写$f=dg$，其中$g$本原。因$d\mid a_n$，$p\nmid d$。素因子指数比较说明$g$仍满足所列三个条件，因此它在$R[x]$中不能分为两个正次数因子。由前一推论，$g$在$F[x]$中不可约，而$f$与$g$在$F[x]$中相伴，所以$f$也不可约。
\end{proof}

\begin{example}
对素数$p$和正整数$n$，$x^n-p$在$\Q[x]$中不可约：$n>1$时应用判别法，$n=1$时为一次多项式。高斯整数中的$\pi=1-\sqrt{-1}$范数为$2$，故不可约；若$a_1,\ldots,a_{n-1}\in(\pi)$，则$x^n+a_{n-1}x^{n-1}+\cdots+a_1x+\pi$同样由判别法在$\Q(\sqrt{-1})[x]$中不可约。
\end{example}

\section{希尔伯特基定理}

\begin{definition}
交换环的每条理想升链都终止时，称其为诺特环。若其理想构成的每个非空集合按包含关系都有极大元，称其满足理想的极大条件。这里的极大元只要求在所考察集合内极大。
\end{definition}

\begin{theorem}\label{p2:noetherian-equivalences}
对交换环$R$，下列条件等价：
\begin{enumerate}
\item $R$为诺特环；
\item $R$满足理想的极大条件；
\item $R$的每个理想均有限生成。
\end{enumerate}
\end{theorem}
\begin{proof}
\begin{enumerate}
\item[$1)\Rightarrow2)$] 若某个非空理想集合$S$没有极大元，先取$N_1\in S$，再逐次选取$S$中严格包含前一项的理想，得到无限严格升链$N_1\subsetneq N_2\subsetneq\cdots$，与诺特性矛盾。
\item[$2)\Rightarrow3)$] 对理想$N$，考虑包含于$N$的全部有限生成理想组成的集合。它含$(0)$，所以有极大元$A$。若$A\ne N$，取$a\in N\setminus A$，则$A+(a)$仍为有限生成理想，且严格包含$A$，矛盾。因此$N=A$有限生成。
\item[$3)\Rightarrow1)$] 给定理想升链$N_i$，其并$N$为理想。写$N=(a_1,\ldots,a_r)$，有限个生成元同属某个$N_m$，故$N\subseteq N_m$。所以$N=N_m=N_{m+1}=\cdots$，升链终止。
\end{enumerate}
\end{proof}

\begin{proposition}
\begin{enumerate}
\item 诺特环的满同态像是诺特环；
\item 若$N$和$R/N$均满足理想升链条件，则$R$也满足理想升链条件；
\item 有限多个诺特环的直和仍为诺特环。
\end{enumerate}
\end{proposition}
% 原文作为习题给出，以下补全。第二条的N视为可能无单位元的环。
\begin{proof}
\begin{enumerate}
\item 对满同态取目标理想链的完全反像，得到源环中的理想链，源环链稳定后，再取像即知目标链稳定。
\item 设$A_1\subseteq A_2\subseteq\cdots$为$R$的理想链。$A_i\cap N$为$N$的理想，$(A_i+N)/N$为$R/N$的理想；两条链从某个共同的$m$起均稳定。对$i\ge m$及$a\in A_i$，由像相等，存在$b\in A_m$使$a-b\in N$。于是$a-b\in A_i\cap N=A_m\cap N$，故$a\in A_m$。所以$A_i=A_m$。
\item 两个环的直和中，以一个分量为理想，商环同构于另一个分量，用第二条即可；有限多个分量归纳得到。
\end{enumerate}
\end{proof}

\begin{theorem}[希尔伯特基定理]\label{p2:hilbert-basis}
若$R$为有单位元的诺特环，则$R[x]$也是诺特环。
\end{theorem}
\begin{proof}
设$H$为$R[x]$的理想，求证它有限生成。对每个$r\ge0$，取$H$中次数至多为$r$的多项式的$x^r$系数全体，记为$I_r$。等价地，它由$0$与$H$中$r$次多项式的首项系数组成。

若$a,b\in I_r$，取相应$f,g\in H$，则$f-g$的$x^r$系数为$a-b$，所以$a-b\in I_r$。对$c\in R$，$cf\in H$给出$ca\in I_r$，即使首项抵消得到$0$也仍在其中。因此$I_r$为理想。由$xf\in H$，还有$I_r\subseteq I_{r+1}$。

诺特性保证存在$m$使$I_m=I_{m+1}=\cdots$。各$I_r$有限生成，写
\[
I_r=(a_{r1},\ldots,a_{rn_r})\qquad(0\le r\le m),
\]
并取$f_{ri}\in H$，其次数为$r$、首项系数为$a_{ri}$。零理想$I_r$不需要选取生成多项式。

求证有限多个$f_{ri}$生成$H$。对$f\in H$按次数$t$归纳，$f=0$时成立。设$f=ax^t+\cdots\ne0$，令$r=\min(t,m)$。因$a\in I_t=I_r$，存在$b_i\in R$使$a=\sum_i b_ia_{ri}$。取
\[
g(x)=\sum_i b_i x^{t-r}f_{ri}(x).
\]
% 高次数情形必须乘x^{t-r}，才能与f的首项次数一致。
则$g\in H$且首项恰为$ax^t$，所以$f_1=f-g\in H$，且$f_1=0$或$\deg f_1<t$。由归纳假设，$f_1$为这些$f_{ri}$的组合，$g$本身也是这样的组合，故$f$也在它们生成的理想中。特别地，$t=0$时相消后直接为零，归纳起点成立。因此$H$有限生成，$R[x]$为诺特环。
\end{proof}

\begin{corollary}
若$R$为交换诺特幺环，则：
\begin{enumerate}
\item $R[x_1,\ldots,x_n]$为诺特环；
\item 任一与$R$共用单位元的有限生成扩环$R[u_1,\ldots,u_n]$为诺特环，且关系理想
\[
I=\{f:f(u_1,\ldots,u_n)=0\}
\]
有限生成。
\end{enumerate}
\end{corollary}
\begin{proof}
\begin{enumerate}
\item 对未定元个数归纳，逐次应用基定理。
\item 代入同态将诺特环$R[x_1,\ldots,x_n]$满射到生成环，所以它为诺特环；核$I$是诺特环的理想，故有限生成。这说明全部代数关系可由有限条基本关系生成。
\end{enumerate}
\end{proof}

\begin{definition}
设$R$为诺特环，$A,B$为理想。若$A\subseteq B$，称$B$为$A$的因子；若严格包含，称为真因子。若$A$不能写成两个真因子的交，称$A$为不可约理想。
\end{definition}

\begin{theorem}\label{p2:noetherian-irreducible-intersection}
诺特环的每个理想均为有限多个不可约理想之交。
\end{theorem}
\begin{proof}
反设不能作这种表示的理想组成非空集合$S$。由极大条件，取其极大元$A$。$A$不能不可约，故$A=B\cap C$，其中$B,C$严格包含$A$。由极大性，$B,C\notin S$，分别可写成有限多个不可约理想之交。将两种表示合并，就得到$A$的有限不可约理想交表示，与$A\in S$矛盾。
\end{proof}

\begin{example}
素理想是不可约理想。事实上，若$P=B\cap C$且$B,C$都严格包含$P$，取$b\in B\setminus P,c\in C\setminus P$。则$bc\in B\cap C=P$，与素性矛盾。

极大理想的幂未必不可约。在$\Z[x]$中，$P=(2,x^2+x+1)$为极大理想，因为商环同构于四元域$\F_2[x]/(x^2+x+1)$。令$h=x^2+x+1$，则
\[
P^2=(4,2h,h^2)=(2,h^2)\cap(4,h).
\]
% 补充验证交等式，并说明它确为真分解。
将$\Z[x]$视为$\Z[h]$模上的自由模，基为$1,x$：用首一二次式$h=x^2+x+1$反复消去$x^2$即可得到该表示，比较最高次项又可得唯一性。逐系数考虑$\Z[h]$中的两个理想，属于$(2,h^2)$要求常数项及一次项系数为偶数，属于$(4,h)$要求常数项被$4$整除。因此交中常数项被$4$整除、一次项被$2$整除，高次项任意，恰为$(4,2h,h^2)$。而$2\notin P^2$、$h\notin P^2$，两个因子均严格大于$P^2$。
\end{example}

\begin{theorem}\label{p2:z-polynomial-primes}
$\Z[x]$的素理想恰分为三类：
\begin{enumerate}
\item $(0)$；
\item 由一个不可约元生成的非零主理想，此类理想不极大；
\item $(p,f(x))$，其中$p$为素数，$f$为正次数整系数多项式，其首项系数与$p$互素，且$\overline f\in\F_p[x]$不可约。这些理想也都是极大理想。
\end{enumerate}
第三类中，$(p,f)=(q,g)$当且仅当$p=q$且$\overline f,\overline g$在$\F_p[x]$中相伴。
\end{theorem}
\begin{proof}
$\Z[x]$是高斯整环，故非零主理想$(f)$为素理想当且仅当$f$不可约。这样的主理想不极大：若$f=p$为整数素数，则商环$\F_p[x]$不是域；若$f$为正次数本原不可约多项式，取一个不整除其首项系数的素数$p$，$\overline f$仍有正次数，因而在$\F_p[x]$中有不可约因子$h$。理想$(p,\widetilde h)$为包含$(f)$的真理想，且严格包含它，因为$p\notin(f)$。这里$\widetilde h$表示$h$的整系数提升。

设$I$为非零非主素理想。在$I$的非零元素中，先取次数最低者，再在该次数下取正首项系数最小者$p_0(x)$。若$I=(p_0)$，与非主性矛盾，所以可取$p_1\in I\setminus(p_0)$，并使其次数最低。

先证$p_0$不可约。若$p_0=fg$，由素性，$f$或$g$属于$I$。若两个因子均为正次数，就与$p_0$次数最低矛盾。因此其中一个为整常数。若两个因子均非单位，则属于$I$的那个因子或者次数更低，或者同次数但正首项系数更小，仍矛盾。所以$p_0$不可约。

再证$\deg p_0=0$。反设其次数为正，则$p_0$本原。令$a$为它的首项系数；通过反复将$p_1$乘以$a$并减去$p_0$的适当单项式倍数，可得到
\[
a^kp_1=qp_0+r,\qquad \deg r<\deg p_0,
\]
其中$q,r\in\Z[x]$。$r\ne0$，否则在$\Q[x]$中$p_0\mid p_1$；由本原性及高斯引理便有$p_0\mid p_1$在$\Z[x]$中成立，矛盾。而$r=a^kp_1-qp_0\in I$又违反$p_0$的选择。因此$p_0=p$为整数素数。

由$(p)\subseteq I$，令$\overline I=I/(p)\subseteq\F_p[x]$。它为非零素理想，故由某个首一不可约多项式$\overline f$生成。选择其整系数提升$f$，由理想对应，
\[
I=(p,f),\qquad
\Z[x]/I\cong\F_p[x]/(\overline f),
\]
右侧为域，所以$I$极大。反过来，所列$(p,f)$的商环均为域，故均为极大理想和素理想。

最后，$(p,f)\cap\Z=(p)$，所以两个这样的理想相等必有$p=q$。固定$p$后，再用模$p$的理想对应及$\F_p[x]$中主理想相等的判别，就得到$\overline f,\overline g$相伴的充要条件。
\end{proof}
